% Appendix F: simulation specification. Referenced by S3 (budget normalization,
% parameter defaults, scoring), and by every simulation figure's caption
% ("Settings: Appendix D"). Carries the generating context removed from captions and
% footnotes in the 2026-09-27 compression. Code: the model repository (commit hashes
% 21b0d9a for the canonical sweeps and d77c854 for the production-function sweeps).
\section{Simulation specification}\label{app:simulation}

\subsection*{Budget and productivity}
Over the whole horizon, the funder's total budget is $b\,n\,\mathbb{E}[R_{0}]$, where $\mathbb{E}[R_0]$ is the expected baseline resources of one researcher, so at $b = 1$ the budget equals the community's expected baseline resources for one round. The budget does not depend on the horizon $T$. Strategies that do not plan ahead spend $1/T$ of the budget each round. For the timing runs of Appendix~\ref{app:timing} (Figures~\ref{fig:depth-schedule} and~\ref{fig:timing-boundary}) the budget is scaled against capability instead, $b\,n\,\mathbb{E}[K_0]$, so that a community with few resources does not zero the budget. The simulations set the productivity constant $A = 1$. The code parameterizes the budget as $2\,b'\,n\,\mathbb{E}[R_0]$; its $b'$ is half of the $b$ reported here.

\subsection*{Default parameters}
Unless a figure's settings below say otherwise: $T = 2$ rounds; $n = 50$ researchers; capability and baseline resources drawn from Pareto distributions with scale 1 and tail parameters $\alpha_K = \alpha_R = 2$; copula parameter $\rho = 0$; compounding rate $\epsilon = 0.1$; budget scale $b = 1$; review noise $\tau_K = 1$ and resource-signal noise $\tau_R = 1$; seed-grant fraction $0.25$ where seed grants are used.

\subsection*{Population draws, particles, and heavy tails}
Capability and resources are coupled through a Gaussian copula: standard normal draws with correlation $\rho$ are mapped through the normal distribution function to uniforms and then through the Pareto quantile functions. A Pearson correlation between two Pareto variables is undefined when either tail parameter is at most 2, so $\rho$ is a copula parameter and not a correlation of the draws; the rank correlation it induces is $(6/\pi)\arcsin(\rho/2)$, which is $0.48$ at $\rho = 0.5$ and $0.79$ at $\rho = 0.8$. Below a tail parameter of 2 the Pareto distribution has infinite variance and sample means converge slowly: at $\alpha_K = 1.3$ with the scale set for a mean of 2, the mean capability across 100 populations of 50 is 1.77 and the median population has mean 1.48, because the rare draws that carry the mean are absent from most samples. Heavy-tailed results are therefore reported with standard errors across populations, and Table~\ref{tab:particles} reports medians beside means.

Bayesian funders compute each researcher's posterior by importance sampling: $M$ particles $(K, R_0)$ are drawn from the prior afresh in each round and weighted by the full likelihood of the researcher's record and signals to date. With few particles the posterior can rest on a handful of them, and rare regions of the prior, where the most capable researchers of a heavy-tailed field sit, can be missed. Table~\ref{tab:particles} checks this. Raising $M$ from 200 to 1{,}000 raises the value of review by 0.01 to 0.02 in every field at every noise level (the posteriors at $M = 200$ understate the capability of the most capable researchers), and raising it further to 5{,}000 changes the value of review by at most 0.005, within one standard error everywhere. The value of targeting uses no particles. Figures~\ref{fig:review-value} and~\ref{fig:records-rounds}, the sensitivity analysis of Appendix~\ref{app:sensitivity}, and Figure~\ref{fig:regime} use $M = 1{,}000$; the remaining simulation figures use $M = 200$, and by the table their values of review are understated by about 0.01 to 0.02 of the no-review shortfall, which does not change any comparison they support.

\begin{table}[t]
\centering
\small
\caption{Convergence in the number of particles. The value of review (share of the no-review shortfall recovered; mean over 100 populations, with the median in parentheses) and the effective sample size of the round-one posteriors (median across researchers and populations, with the tenth percentile in parentheses), by field, review noise, and number of particles $M$. Two rounds, $b = 1$, mean capability 2; the same 100 populations at every $M$. The value of targeting uses no particles and is unchanged.}
\label{tab:particles}
\begin{tabular}{@{}llccc@{}}
\toprule
$\alpha_K$ & $\tau_K$ & $M = 200$ & $M = 1000$ & $M = 5000$ \\
\midrule
\multicolumn{5}{@{}l}{\emph{Value of review, mean (median)}} \\
1.3 & 0.3 & 0.91 (0.94) & 0.94 (0.96) & 0.94 (0.97) \\
1.3 & 1 & 0.77 (0.83) & 0.79 (0.84) & 0.79 (0.85) \\
1.3 & 3 & 0.48 (0.56) & 0.50 (0.58) & 0.50 (0.53) \\
2 & 0.3 & 0.84 (0.87) & 0.86 (0.89) & 0.87 (0.89) \\
2 & 1 & 0.63 (0.67) & 0.65 (0.69) & 0.65 (0.69) \\
2 & 3 & 0.28 (0.26) & 0.30 (0.26) & 0.29 (0.27) \\
3.5 & 0.3 & 0.62 (0.65) & 0.63 (0.66) & 0.64 (0.67) \\
3.5 & 1 & 0.31 (0.32) & 0.32 (0.33) & 0.32 (0.33) \\
3.5 & 3 & 0.05 (0.04) & 0.06 (0.05) & 0.06 (0.04) \\
\addlinespace
\multicolumn{5}{@{}l}{\emph{Effective sample size, median (tenth percentile)}} \\
1.3 & 0.3 & 81 (5) & 396 (28) & 1998 (136) \\
1.3 & 1 & 122 (15) & 618 (73) & 3081 (369) \\
1.3 & 3 & 142 (36) & 711 (164) & 3559 (796) \\
2 & 0.3 & 74 (6) & 368 (33) & 1843 (166) \\
2 & 1 & 124 (18) & 616 (90) & 3079 (447) \\
2 & 3 & 146 (46) & 732 (227) & 3667 (1137) \\
3.5 & 0.3 & 92 (12) & 460 (56) & 2300 (294) \\
3.5 & 1 & 136 (39) & 685 (193) & 3417 (955) \\
3.5 & 3 & 153 (67) & 771 (314) & 3857 (1558) \\
\bottomrule
\end{tabular}
\end{table}

\subsection*{Scoring and comparison}
Each strategy is scored by the sum of expected research output along its run, which equals realized research output averaged over repetitions of the same allocations; this removes simulation noise from comparisons without favoring any strategy. Every configuration is run on independently drawn populations, with the same random draws (populations, signals, and realized research output) used for every strategy, so that comparisons between strategies are paired. Results are reported as shares of a named comparison, stated in each figure: most often the research output that funding adds (a strategy's expected research output minus the expected research output under no funding), and otherwise the no-review shortfall or the optimal allocation's gain. Standard errors, where shown, are across simulated populations.

\subsection*{Settings by figure}
\begin{table}[h]
\centering
\small
\begin{tabular}{@{}p{2.1cm}p{10.2cm}@{}}
\toprule
Figure & Settings that differ from the defaults \\
\midrule
\ref{fig:review-value} & Review noise $\tau_K$ from 0.05 to 20; $\alpha_K \in \{1.3, 2, 3.5\}$ with the Pareto scale set to $2(\alpha_K - 1)/\alpha_K$ so that mean capability is 2 in every field; 100 populations per point; $M = 1{,}000$. The no-review shortfall relative to complete information is 44, 26, and 11 percent of the research output that complete-information funding adds, at the three tail parameters. \\
\ref{fig:records-rounds} & $T = 20$; 50 populations; $M = 1{,}000$. No-review is the myopic funder without review; with review is the myopic funder with a single review score observed at the start. The shortfall is the funder's expected research output below the complete-information optimum in that round, as a share of the optimum's gain over no funding. Under heavy-tailed capability ($\alpha_K = 1.3$, mean 2) the no-review shortfall falls from 60 to 14 percent against 6 to 2 percent with review, and the grant correlations run 0.12 to 0.55 against 0.89 in round one. \\
\ref{fig:overtrust} & True noise $\tau_K = 3$; believed noise from 0.1 to 10; $\alpha_K \in \{1.3, 2\}$; 200 populations. Each negative value is within two standard errors of zero. \\
\ref{fig:depth-schedule} & $T = 6$; community with few resources (Pareto scale for resources 0.001); $\alpha_K = 1.3$; no review signal ($\tau_K = 100$); $\epsilon = 10^{-4}$; budget scaled against capability; $b = 1$ (50 populations) and $b = 6$ (24 populations, standard error of each share at most 0.006). \\
\ref{fig:timing-boundary} & $T = 5$; forward-looking funder with review; budget scaled against capability; resource scales 0.001, 0.3, and 3; $\epsilon$ from $10^{-4}$ to 0.85; 200 populations. Five resource scales were run; the figure shows three, and the other two lie between them. \\
\ref{fig:schedule-vs-signal} & $T$ from 2 to 10 (to 5 for $\epsilon \in \{0.05, 0.15, 0.55\}$); $\epsilon \in \{0.05, 0.15, 0.3, 0.55, 0.85\}$; 64 to 200 populations per point. The gain is the forward-looking funder's research output over the myopic funder's, both with review; negative gains, all within noise of zero, are shown as zero. \\
\ref{fig:floor-cost} & Review-informed forward-looking funder; seed-grant fraction from 0 to 0.75; heavy-tailed field ($\alpha_K = 1.3$) with reliable review ($\tau_K = 0.3$) at $b \in \{0.2, 1, 2\}$, and the intermediate field at $b = 1$; 200 populations. The research output lost is reported as a share of the review-informed funder's gain over no funding in the same setting. \\
\ref{fig:lottery-schemes} & Single round; $n = 50$; a fifth of the population funded; review score $K_i$ plus Gaussian noise of standard deviation $\tau$; expected research output computed exactly given each allocation; 2{,}000 populations per point (Monte Carlo standard error at most 0.002). Left: $\alpha_K = 1.3$, $b = 0.2$. Right: $\alpha_K = 3.5$, $b = 1$. \\
\ref{fig:signal-robustness}, Table~\ref{tab:timing-technology} & Production functions with $\gamma \in \{0, -3\}$ and Leontief; 100 populations per point (50 for Leontief); $T = 2$ for the signal value and $T = 5$ for the timing table. These runs use the greedy allocator, the only one implemented for the general family (the harmonic-mean runs use exact water-filling), so their magnitudes are compared with one another and not with the body's figures. Leontief allocations use a finer allocation step, with ties at the kink broken by fill order. \\
\ref{fig:concentration} & Single round; review-informed funder; $\gamma$ from 0 to $-12$ and Leontief; three field-and-budget combinations ($\alpha_K = 1.3$, $b = 0.2$; $\alpha_K = 2$, $b = 2$; $\alpha_K = 3.5$, $b = 2$); 200 populations per point, 400 at the refinement points of the evenly spread curve. \\
\ref{fig:gap-gini}, Tables~\ref{tab:sens-mean} to~\ref{tab:sens-contam} & $T = 2$; $b = 1$; myopic funders with and without review and the complete-information benchmark; $\rho \in \{-0.5, 0, 0.5, 0.8\}$; $\alpha_R \in \{1.3, 2, 3.5\}$; $n \in \{25, 50, 100, 200\}$; Pareto scale 1 or set to give mean capability 2; review score $K_i + \beta R_{i0}$ plus noise with $\beta \in \{0, 0.25, 0.5, 1\}$; resource-signal noise $\tau_R \in \{0.1, 0.3, 1, 3, 10\}$ and no resource signal; multiplicative review noise $\log S = \log K + e$ with $e$ Gaussian of standard deviation 0.1 to 1.5; review redrawn every round ($T = 2$, and $T = 20$ for the by-round comparison); 50 populations per setting, 100 in the noise sweeps; $M = 1{,}000$. \\
\ref{fig:regime} & $T = 2$; $b \in \{0.01, 0.02, 0.05, 0.1, 0.2, 0.5, 1, 2, 3\}$; $\alpha_K \in \{1.2, 1.3, 1.5, 2, 2.5, 3.5, 5\}$ with mean capability 2; $\tau_K = 1$; myopic funders with and without review and the complete-information benchmark; 50 populations per cell; $M = 1{,}000$; contours interpolated linearly in $\log b$ and in the Gini coefficient of the capability distribution, $1/(2\alpha_K - 1)$. \\
Table~\ref{tab:dynopt} & Complete information; $T = 2$ exact optimum and $T = 5$ optimum over the division of the budget across rounds; $\alpha_K \in \{1.3, 2, 3.5\}$ with mean capability 2; $\epsilon \in \{0.1, 0.5, 0.85\}$; $b \in \{0.2, 1, 2\}$; 20 populations per setting. \\
\bottomrule
\end{tabular}
\end{table}

\subsection*{Placing programs on the budget axis}
A round of the model corresponds to a year of a program's awards, and the regime map's budget $b$ is the total over its two rounds, so a program whose annual awards to a field are a fraction $r$ of the field's annual research spending from every other source, which stands in for the baseline resources of the field's researchers, sits at $b = 2r$. The estimate counts money and not time, and treats each funder's field as the disciplines it mainly funds, so it is an order of magnitude; the bands in the figure span a factor of about two around each placement. United States universities spent \$108.8 billion on research and development in fiscal year 2023, \$55.2 billion of it in the life sciences, and received \$33.1 billion from the Department of Health and Human Services, chiefly the National Institutes of Health, and \$6.7 billion from the National Science Foundation \citep{NCSES2025HERD}: a biomedical agency of that size has $r$ near 1.5 (\$33 billion against \$22 billion from other sources), so $b$ near 3, at the top of the map; a science agency of that size funding the other half of academic research has $r$ near 0.15 and $b$ near 0.3. The European Research Council's budget of 16 billion euros over 2021 to 2027 \citep{ERC2025}, about 2.3 billion a year, against 86 billion euros of higher-education research spending in the European Union in 2024 \citep{Eurostat2025}, gives $r$ near 0.03 and $b$ near 0.05. A private foundation awarding on the order of \$100 million a year across several fields has $r$ at or below 0.01 in any one of them, and $b$ at or below 0.02.

\subsection*{Code}
The model, the sweeps, and the scripts that produce every figure and number in the paper are in the project repository; the archived version will be cited at publication. % [check: repository URL and archived commit at publication]
